\section{引言}
\label{sec:intro}
%------------------------------------------------------------------------------
质子与中子合称核子，是当今宇宙中可见物质的基本构件。自约一个世纪前在实验室被发现以来~\cite{Rutherford:1919fnt,Chadwick:1932ma}，它们的基本性质得到了深入的探索：从通过液氢比热测定自旋~\cite{Dennison}，到磁矩的测量~\cite{Stern}，再到通过弹性电子散射提取其电磁尺寸~\cite{Hofstadter:1956qs}。然而最具揭示性的发现来自 1960 年代末在斯坦福直线加速器中心（SLAC）对质子与原子核所做的电子深度非弹性散射（DIS）实验，质子与中子的组分——夸克（以及后来的胶子）——由此被发现~\cite{Bloom:1969kc}。此后不久，建立在``色''SU(3) 规范对称性基础上的量子场论（QFT）——量子色动力学（QCD）——被确立为强相互作用的基本理论~\cite{Fritzsch:1973pi,Gross:1973id,Politzer:1973fx}，也是核子内部结构的基本理论~\cite{Thomas:2001kw}。

在过去的五十年里，无论在实验还是理论上，对核子内部结构的理解都取得了重大进展。人们建造了多个实验装置来研究涉及质子与原子核的高能碰撞，积累了大量实验数据。基于超越 Feynman 部分子模型~\cite{Feynman:1973xc}的微扰 QCD 分析所导出的 QCD 因子化定理~\cite{Collins:2011zzd}，表征以无穷大动量运动的强子中夸克与胶子纵向动量分布的部分子分布函数（PDF）已通过对这些数据的全局拟合得到~\cite{Gao:2017yyd, Hou:2019efy, Harland-Lang:2014zoa, Ball:2017nwa}。近期一个唯象质子 PDF 的结果示于 Fig.~\ref{fig:PDFS}，其中 $x$ 是部分子所携带的质子动量份额。PDF 对核子中夸克与胶子的内容提供了全面的描述。在理论前沿，QCD 的欧氏路径积分表述结合格点正规化与蒙特卡罗模拟~\cite{Wilson:1974sk}，为强相互作用的非微扰计算提供了一种系统地进行 {\it ab initio}（从头）计算的途径。计算能力的迅速提升与智能数值算法的发展，使格点 QCD 方法在计算强子谱、强耦合常数、强子形状因子乃至散射相移等方面取得了极大成功~\cite{Tanabashi:2018oca,Briceno:2017max,Aoki:2019cca}。

\begin{figure}[htb]
	\includegraphics[width=0.9\linewidth]{images/intro-PDFCT18_fig01.pdf}
	\caption{CTEQ-TEA 合作组（CT18）通过对全局高能散射数据拟合得到的唯象部分子分布~\cite{Hou:2019efy}，其中 $0\le x\le 1$ 为部分子携带的质子无穷大动量的份额。}
	\label{fig:PDFS}
\end{figure}

尽管有这些令人瞩目的成就，我们仍未能从第一性原理系统地解释质子的部分子结构；更明确地说，我们在从 QCD 拉氏量出发计算夸克与胶子分布方面尚未取得根本性进展（简要总结见 \sec{intro_other}）。这背后其实有一个很好的理由：教科书~\cite{Sterman:1994ce,Collins:2011zzd}中部分子物理的标准表述是通过定义在光前（LF，$t-z={\rm const.}$）上夸克与胶子场的{\it 动力学关联函数}完成的，其重要特征是与质子的动量无关。另一方面，格点 QCD 在虚时间的欧氏空间中表述，无法直接计算依赖实时间的动力学关联。处理部分子物理的标准格点方法是计算部分子分布的低阶矩，即局域算符的矩阵元~\cite{Lin:2017snn}。然而，局限于最初几阶矩使研究者无法可靠地重现 Fig.~\ref{fig:PDFS} 所示依赖于 $x$ 的结构，而只能拟合模型的函数形式。多年来，人们提出了光前量子化（LFQ）中的哈密顿量对角化~\cite{Brodsky:1997de}与 Schwinger-Dyson 方程~\cite{Maris:2003vk}等 Minkowski 途径来求解核子结构。虽然在唯象上有显著进展，但计算核子 PDF 的系统近似方法仍然缺失。

几年前，本文的部分作者基于 Feynman 关于部分子的原始思想提出了一种计算依赖于 $x$ 的部分子分布的一般途径：部分子是{\it 大动量下质子静态性质}的无穷大动量极限，因而是本质上可通过格点 QCD 计算的欧氏量~\cite{Ji:2013dva,Ji:2013fga,Ji:2014gla}。据此，中间区域 $x_{\rm min}\sim 0.1 <x<x_{\rm max}\sim 0.9$ 的部分子物理可以从具有中等偏大动量（例如 Lorentz 推进因子 $\gamma=2-5$）的质子的物理性质计算得到。该理论被命名为{\it 大动量有效理论}（LaMET），因为在无穷大动量系（IMF）的部分子与有限动量的夸克、胶子之间建立严格联系，需要在有效场论（EFT）中对大动量的紫外（UV）模作恰当处理并进行系统的幂次计数。

LaMET 的基本原理来自朴素部分子模型中一个隐含的观察：只要质子动量远大于典型的强相互作用标度 $\Lambda_{\rm QCD}$ 或其质量，其结构就近似地与动量无关。例如，动量为 $P=|\vec{P}|=5$ GeV 的质子在中等 $x$ 处的夸克动量分布与 $P=50$ GeV 或 $P=5$ TeV 时没有太大差别。人们可以把这一现象称为{\it 大动量对称性}，其性质类似于氢原子的电子结构对质子质量不敏感——只要后者远大于电子质量。质子结构的渐近行为或许由 $\Lambda_{\rm QCD}/P$ 展开控制，但要给出论证需要对底层动力学有更好的理解。在此假设下，Feynman 用无穷大动量 $P=\infty$ 的质子代替高能散射中被探测的大而有限的动量的质子，对应于 $\Lambda_{\rm QCD}/P$ 展开的首项，于是{\rm IMF 中的质子}及其组分——{\it 部分子}——这些理想化概念便诞生了。

然而在 QFT 中，$P=\infty$ 极限的存在性取决于理论的 UV 行为。一般而言，无穷大动量极限与 UV 截断极限 $\Lambda_{\rm UV}\to \infty$ 并不对易。物理极限是 $(\Lambda_{\rm UV}\gg P) \to\infty$，而部分子模型及随后的 QCD 因子化定理使用 $(P\gg\Lambda_{\rm UV})\to\infty$，保持所有 PDF 具有有限支撑 $|x|\le 1$，其中负 $x$ 对应反夸克。因此部分子是一种在真实世界中并不存在的理想化概念。幸运的是，由于渐近自由，上述差别可以在微扰 QCD 中计算出来。因此，LaMET 是关于部分子的有效理论：它使用通常的场论计算次序 $(\Lambda_{\rm UV}\gg P)\to\infty$，通过 EFT 匹配与演化系统地处理不对易的 $P\to\infty$ 极限，并以幂次修正处理有限 $P$ 效应。这样，定义在 IMF 或 LF 上的 PDF 可以在中等 $x$ 区域由 $P\sim $ 几 GeV 的结构计算获得。

LaMET 的首个应用是极化质子中的总胶子螺旋度 $\Delta G$，这一物理量在极化 RHIC 上具有重要的实验意义~\cite{Bunce:2000uv}，却多年来超出理论可及范围。文献~\cite{Ji:2013fga}中我们证明，从物理规范下胶子自旋算符的大动量矩阵元出发，$\Delta G$ 可以通过 EFT 匹配得到。继此成功之后，LaMET 被应用于共线夸克 PDF~\cite{Ji:2013dva}。后一应用引发了可观的理论与数值研究活动，尤其是针对质子及其他强子中味单态外的 $u-d$ 分布。随后文献~\cite{Ji:2014gla}建立了一般的 LaMET 框架。最近，该方法也被推广到胶子情形~\cite{Zhang:2018diq,Li:2018tpe}。因此，如今可以直接在特定 Feynman 变量 $x$ 处用格点 QCD 计算 PDF，而无需求助于 LFQ。此外，质子的部分子图景极为丰富，LaMET 有望计算共线 PDF 之外的部分子物理。

近年来，为质子定义新的部分子观测量取得了巨大进展。特别是，在刻画质子的横向结构时发展出两个并行概念。第一个是广义部分子分布（GPD）~\cite{Mueller:1998fv,Ji:1996ek,Radyushkin:1998es}。GPD 综合了质子弹性形状因子（给出部分子在横向空间的密度~\cite{Miller:2007uy}）与 Feynman PDF 的特点并在二者之间插值。给定纵向动量与横向空间的联合分布，人们可以构造部分子的轨道角动量（OAM）等物理量~\cite{Ji:1996ek}。一般地，GPD 可用于生成按动量剖析的质子横向空间图像~\cite{Burkardt:2000za}。一类新的实验过程——深度虚 exclusive 过程（DVEP），包括终态为衍射实光子加反冲质子的深度虚 Compton 散射（DVCS）——被发现可以测量 GPD~\cite{Ji:1996ek,Ji:1996nm}。第二个概念是横动量依赖（TMD）PDF（或 TMDPDF），其中部分子的横动量被显式地保留~\cite{Collins:1981uk,Collins:2011zzd}。近年来在其恰当定义、因子化与自旋关联方面取得了大量理论进展~\cite{Collins:2012uy,Echevarria:2012js,Collins:2017oxh}。TMDPDF 可以通过观测终态粒子的横动量在实验过程中测量。

\begin{figure}[htb]
	\centering
	\includegraphics[width=0.95\linewidth]{images/intro-EIC_fig02.pdf}
	\caption{布鲁克海文国家实验室（BNL）电子离子对撞机的一个设计方案（图片由 BNL 提供），可用于探测质子的部分子图景。}
	\label{fig:EIC.png}
\end{figure}

多年来人们逐渐认识到，需要一个专用实验装置来全面探索质子的部分子图景。为满足这一需求，美国核科学界提议建造高能量、高亮度电子离子对撞机（EIC）~\cite{Geesaman:2015fha}，并已于近期获美国能源部批准。这台新对撞机将电子加速到 10--30 GeV，将离子——包括质子以及直至 Pb 或 U 的重核——加速到每核子最高 100 GeV，实现 40 至 170 GeV 的质心碰撞能量 $E_{\rm cm}$。对应的固定靶实验等效电子能量为 100 GeV 至 10 TeV。束流是极化的，亮度高达 $10^{33-34}$ 碰撞/(${\rm cm}^2\cdot{\rm s}$)，这对研究 DVCS 等 exclusive 过程至关重要。碰撞的运动学范围覆盖低至 $10^{-4}$ 以下的 Bjorken $x_B$（在下一节讨论的朴素部分子模型中它与部分子动量份额 $x$ 一致），以及高达 $10^4$ GeV$^2$ 的 $Q^2$。EIC 的科学目标已在专门研究中详述~\cite{Accardi:2012qut}。

当然，EIC 与格点 QCD 的努力不会止步于质子的精密部分子物理。更重要的是，我们需要发展描述作为强耦合相对论量子系统的核子的方式或语言，就像凝聚态物理中对量子霍尔效应的理解那样。若不深入理解强耦合 QCD 物理的机制，我们就不能宣称从根本上理解了质子与中子的结构，特别是其质量与自旋的起源。这是当今粒子与核物理标准模型面临的最具挑战性的目标之一。

本综述旨在系统阐述 LaMET 处理部分子物理的思想、形式体系与结果。由于该领域发展迅速，我们不追求完全完备；相关领域的参考文献也不求全，对可能的遗漏深表歉意。与格点部分子物理密切相关的综述可见~\cite{Cichy:2018mum,Zhao:2018fyu}。已有若干工作检验 LaMET 在各种模型中的有效性~\cite{Ji:2018waw,Gamberg:2014zwa,Broniowski:2017wbr,Xu:2018eii,
	Son:2019ghf,Ma:2019agv,Nam:2017gzm,Bhattacharya:2018zxi,Jia:2015pxx,Hobbs:2017xtq,Broniowski:2017gfp,
	Radyushkin:2017gjd,Bhattacharya:2019cme,Kock:2020frx,DelDebbio:2020cbz}，下面我们会择要提及以作说明。也有论文质疑 LaMET 方法的有效性~\cite{Carlson:2017gpk,Rossi:2017muf,Rossi:2018zkn}，其中一些疑问后来已在文献中得到澄清~\cite{Briceno:2017cpo,Ji:2017rah,Radyushkin:2018nbf}；此处不再讨论，有兴趣的读者可参阅上述文献。正文中多数场合使用{\it 质子}一词，以强调其在核物理与粒子物理中的重要性；但相关讨论同样适用于中子及其他强子。

本综述的结构如下。在本引言的剩余部分，我们将阐明部分子物理的本质——它是大动量下质子内部结构的有效描述——以及文献中求解部分子结构的其他现有方法。\sec{lamet} 从运动强子中物理可观测量的动量重正化群方程（RGE）出发引入 LaMET 方法，随后给出动量分布与 PDF 之间的匹配；然后我们表述一个 EFT 展开，用以从适用于大动量质子结构的理论方法计算部分子物理。\sec{renorm} 讨论共线 PDF 的若干重要细节：非局域算符的重正化（尤其是格点正规化中的幂次发散），以及全阶微扰下的匹配。\sec{gpo} 致力于推广到一般的共线部分子观测量，包括 GPD、部分子分布振幅以及更高阶的部分子关联；我们还讨论其在极化质子部分子 OAM 中的应用。\sec{tmd} 考虑在一类新的部分子观测量 TMDPDF 上的应用：研究准 TMDPDF 与物理 TMDPDF 的匹配，并探讨软函数的格点计算。最后，\sec{lattice} 总结与 LaMET 应用相关的近期格点计算，结论见 \sec{conclusion}。综述末尾附有缩略语与术语表以及记号与约定的附录。

%------------------------------------------------------------------------------
\subsection{经由无穷大动量态理解部分子}
\label{sec:lc}
%------------------------------------------------------------------------------
尽管部分子已成为高能散射中无处不在的语言，但它们作为描述核子内部结构的 QCD 有效自由度这一角色在文献中较少被强调。在 QCD 因子化定理的应用中，它们如 Feynman 所设想的那样，是来自无穷大动量极限的对象，潜在的 UV 发散在该极限之后再行正规化和重正化。因此，部分子是一个理想化概念，只有在 IMF 和光前规范 $A^+=(A^0+A^z)/\sqrt 2=0$ 的语境下才指称核子或其他强子的夸克、胶子 Fock 成分。它们与重夸克有效理论（HQET）中的{\it 无穷重夸克}属于同一类概念~\cite{Manohar:2000dt}。为了引出 LaMET，理解部分子的这种来源和本性十分重要。

Feynman 基于非相对论系统（原子与分子）中电子散射的知识~\cite{West:1974ua}，引入了{\it 朴素部分子模型}来描述质子上的深度非弹性散射（DIS），并解释观测到的 Bjorken 标度现象~\cite{Feynman:1969ej,Bjorken:1969ja,Feynman:1973xc}。

\begin{figure}[htb]
	\includegraphics[width=0.6\linewidth]{images/intro-dis_fig03.pdf}
	\caption{深度非弹性散射：质子中的部分子被虚光子探测。}
	\label{fig:DIS}
\end{figure}

Fig. \ref{fig:DIS} 所示为 DIS 过程：动量为 $q^\mu$ 的大动量虚光子被动量为 $P^\mu$、质量为 $M$ 的质子吸收。不变变量为 $Q^2 = -q^\mu q_\mu$ 与 $P\cdot q = M\nu$，且 Bjorken $x_B=Q^2/(2P\cdot q)$ 在标度（或 Bjorken）极限 $Q^2\to \infty$、$P\cdot q\to \infty$ 下固定。 inclusively 的 DIS 截面可以分解为轻子张量与强子张量的乘积，前者与轻子的电磁流相联系，后者包含靶质子电磁相互作用的全部信息。

要了解质子结构，最好在 Breit 系中考察散射：
\begin{align}\label{eq:diskin}
	q^\mu &= (0,0,0,-Q)  \,, \nonumber \\
	P^\mu &=\left(\sqrt{\frac{Q^2}{4x_B^2}+M^2}, 0, 0, \frac{Q}{2x_B}\right) \,,
\end{align}
此时虚光子能量为零。探测只对空间结构敏感，如同非相对论的电子散射。但相对论要求质子以大动量 $P^z= Q/(2x_B)$ 运动，推进因子为 $\gamma = Q/(2x_BM)$，在 Bjorken 极限下趋于 $P^z=\infty$。

Feynman 不纠缠 QFT 的微妙之处，对质子结构与散射机制作出了直观假设~\cite{Feynman:1973xc}：不同大 $P^z$ 下质子的结构应当相似，并可用 $P^z=\infty$（即 IMF）下的结构近似。组分（部分子）之间的相互作用被无限时间膨胀，波函数位形被冻结。高能散射中的质子可视为由无相互作用的部分子组成，每个部分子带有纵向动量 $xP^z$，其中 $0< x< 1$。

非相对论系统的内部结构与整体动量无关；而相对论系统不同，因为它们最少受到 Lorentz 收缩的影响。这类系统的结构与整体运动不可分割地纠缠在一起，其对外部动量的依赖是一个动力学问题。另一方面，如果内部结构依赖于特定的强子标度 $\Lambda_{\rm QCD}$，那么所有满足 $P^z\gg \Lambda_{\rm QCD}$ 的大动量质子都具有相似的结构，对应于 $P^z\to \infty$ 极限。这意味着，若 $f(k^z,P^z)$ 是动量为 $P^z$ 的质子中组分动量 $k^z$ 的分布，它在 $P^z=\infty$ 处可能是解析的，并可展开成 $1/P^z$ 的 Taylor 级数，
\begin{align}\label{Eq:expansion}
	f(k^z, P^z) = f(x) + f_2(x)(\Lambda_{\rm QCD}/P^z)^2 + ...\,,
\end{align}
其中 $x=k^z/P^z$。若是如此，质子性质便存在精确到幂次修正 ${\cal O}(1/P^z)$ 的``大动量对称性''（有时为简单起见我们省略上标 $z$），而 $f(x)$ 即部分子分布。

上述图景在某些简单的 QFT 模型中成立，在这些模型里复合系统波函数对动力学参考系的依赖可以被直接研究。二维系统有许多有趣的例子可以得到解析解。被大量研究的案例之一是大 $N_c$ QCD，也称 't Hooft 模型~\cite{tHooft:1974pnl}，其束缚态具有良定义的大动量极限。波函数可展开为 $1/P$ 的级数，修正从 $(1/P)^2$ 开始。在该极限下组分的动量 $k$ 与 $P-k$ 按比例缩放。当以 $x=k/P$ 为变量作图时，波函数随动量大小的变化可在~\cite{Jia:2017uul} 的 Figs. 8--11 中看到。这正是 Feynman 直觉适用的例子类型。

然而，这样的直觉在许多 3+1 维 QFT（如 QCD）中失效。当束缚态以越来越大的动量运动时，需要越来越多的场论高动量模式来构建其内部结构。Lorentz 收缩意味着对结构重要的组分动量范围也随之增大。如果这些高动量模式不能有效地与低动量模式退耦，结构量中就会出现形如 $\ln P$ 的大对数。因此在这类理论中 $P=\infty$ 处可能存在奇点（割缝），使 $P\to \infty$ 极限病态、大动量展开无法进行。这种情况与理论的 UV 性质密切相关：取 UV 截断 $\Lambda_{\rm UV}\to \infty$ 与 $P\to \infty$ 这两个极限并不对易。物理相关的极限是 $(\Lambda_{\rm UV}\gg P)\to \infty$，而 QCD 因子化中的部分子是在忽略 UV 发散的另一极限 $(\Lambda_{\rm UV}\ll P)\to \infty$ 中获得的。于是可以形式地把部分子分布写成
\begin{align}\label{eq:pd}
	f(x) = \int \frac{d\lambda}{2\pi} e^{ix\lambda}\langle P^z=\infty|\psi^\dagger(z) \psi(0)|P^z=\infty\rangle\,,
\end{align}
其中 $\lambda = \lim_{P^z\to\infty, z\to 0} (zP^z)$，$\psi$ 是量子场。

历史上，场论的 IMF 极限首先在微扰论图形规则层面被研究~\cite{Weinberg:1966jm}。人们发现，忽略 UV 发散而取 $P\to\infty$ 会大大简化微扰论规则：许多按时序图形消失，只有少数有有限贡献。而且，该极限下的散射类似非相对论量子力学，波函数描述变得有用。Fock 态定义了具有正确运动学支撑（$0<x<1$）的部分子。取极限后所有物理量都与 $P$ 无关，在 UV 发散未被正规化之前大动量对称性是严格的。因此，{\it 与 Feynman 朴素部分子模型对应的正是这个``朴素''极限：$\Lambda_{\rm UV}\ll P\to \infty$}。

在高能散射的标准 QCD 研究中，上述作为有效自由度的部分子概念一直被隐式使用。PDF 以朴素的 $P=\infty$ 极限定义，并被用来拟合实验截面，从而形成 QCD 因子化定理~\cite{Collins:2011zzd}。

\subsection{经由光前关联函数理解部分子}
\label{sec:lc-imf}

在文献和教科书中，部分子分布传统上并不像 \eq{pd} 那样用欧氏矩阵元表示，而是用量子场的所谓光前（LF）关联函数（``算符形式''）表示~\cite{Brodsky:1997de,Collins:2011zzd}。用共线量子场和有效拉氏量给出的更显式表述见于软共线有效理论（SCET）~\cite{Bauer:2000yr,Bauer:2001ct,Bauer:2001yt}。

有一个物理的方式可以看出高能散射的部分子描述导致光前关联。考虑质子静止系中的 DIS，其中虚光子动量为
\begin{align}
	q^\mu = (\nu, 0, 0,  -\sqrt{\nu^2+2x_B M\nu}) \,.
\end{align}
在 Bjorken 极限 $\nu\to\infty$ 下，虽然光子的不变质量 $Q$ 趋于无穷，但光子动量实际上趋于类光，即趋近光前。因此，在与光子动量互为 Fourier 共轭的强子张量中，两个电磁流的间隔也趋于光锥方向。

这样看来自然的是：IMF 中质子的全部结构物理也可以用依赖时间的 LF 关联函数或 LF 上量子场的关联来表达。形式上看这一点很简单，只需写出
\begin{align}
	|P\to\infty\rangle = U(\Lambda_\infty) |P=0\rangle \,.
\end{align}
把推进算符 $U(\Lambda_\infty)$ 作用到普通动量分布中的静态非局域算符上，一切静态关联都变成光锥关联。这一推进过程类似于量子力学中哈密顿演化的 Schr\"odinger 图像向 Heisenberg 图像的转变：时间依赖转移到算符上。

为表达光锥关联，方便的做法是引入一对共轭类光（光锥）矢量 $p^\mu=(\Lambda,0,0,\Lambda)$ 和 $n^\mu=(1/2\Lambda,0,0,-1/2\Lambda)$，满足 $n^2=p^2=0$、$n\cdot p=1$，其中 $\Lambda$ 为参数。任何四矢量都可展开为
\begin{align}
	k^\mu = k\cdot n p^\mu + k\cdot p n^\mu  + k_\perp^\mu \ .
\end{align}
特别地，沿 $z$ 方向运动的质子动量 $P^\mu$ 可写为
\begin{align}
	P^\mu = p^\mu + (M^2/2)\, n^\mu \,,
\end{align}
其中 $M$ 是质子质量。

利用上述记号，质子内的非极化夸克分布可表示为~\cite{Collins:2011zzd}
\begin{align}
	q(x) =\frac{1}{2} \int \frac{d\lambda}{2\pi} e^{i\lambda x}
	\langle P|\overline{\psi}(0) \slashed{n} W(0,\lambda n)\psi(\lambda n)|P\rangle_c \ ,
	\label{eq:LCqPDF}
\end{align}
其中 $\psi$ 是夸克场，$W$ 是规范链
\begin{align}
	&W(x_2, x_1)= \nonumber\\
	&\hspace{0.5em}{\cal P}\exp[-ig\int_0^1 dt\,(x_2-x_1)_\mu A^\mu(x_1+(x_2-x_1)t)]
\end{align}
以保证规范不变性，${\cal P}$ 表示路径排序。$c$ 表示仅取连通贡献，后文将其省略。这是规范理论的一个特性：荷场并非规范不变，物理分布必须包含一束共线规范粒子。注意上式对任意动量 $P$ 都成立（一种残余的动量对称性），特别是在核子静止系中也成立。上述分布的 $x$ 支撑为 $[-1,1]$；对负 $x$，定义反夸克分布 $-q(-x)\equiv \bar q(x)$。这一表达式在文献中比 Feynman 最初的 PDF 表述更为人熟知。对单夸克靶，有 $q(x)=\delta(x-1)$。

要在上式中显现出部分子，可以循着 Dirac 于 1949 年建议~\cite{Dirac:1949cp}的 QCD 光前量子化~\cite{Chang:1968bh,Kogut:1969xa,Drell:1970yt}进行。在 LFQ~\cite{Brodsky:1997de}中定义 LF 坐标
\begin{align}
	\xi^\pm = (\xi^0\pm \xi^3)/\sqrt{2} \ ,
\end{align}
其中 $\xi^+$ 是 LF ``时间''，$\xi^-$ 是 LF ``空间坐标''；任何四矢量 $A^\mu$ 写成 $(A^+,A^-,\vec{A}_\perp)$。动力学自由度定义在 $\xi^+=0$ 平面上，具有任意 $\xi^-$ 和 $\vec{\xi}_\perp$，其共轭动量为 $k^+$ 与 $\vec{k}_\perp$。动力学由光锥哈密顿量 $H_{\rm LC}=P^-$ 生成。对于三动量为 $(k^+, \vec{k}_\perp)$、质量为 $m$ 的自由粒子，{on-shell} LF 能量为 $k^-=(\vec{k}_\perp^2+m^2)/(2k^+)$。

对 QCD，可定义 Dirac 矩阵 $\gamma^\pm = (\gamma^0\pm\gamma^3)/\sqrt{2}$ 以及夸克场的投影算符 $P_\pm = (1/2)\gamma^\mp\gamma^\pm$，于是任何 $\psi$ 都可分解为 $\psi=\psi_++\psi_-$，其中 $\psi_\pm = P_\pm \psi$，$\psi_+$ 被视为动力学自由度。对规范场，$A^+$ 被 LF 规范 $A^+=0$ 固定，$A_\perp$ 是动力学自由度；$\psi_-$ 与 $A^-$ 是依赖变量，可用运动方程以 $\psi_+$ 和 $A_\perp$ 表达~\cite{Kogut:1969xa}。

引入正则展开后，LF 关联的物理内涵便一目了然：
\begin{align}
	& \psi_+(\xi^+=0,\xi^-,\vec{\xi}_\perp)
	= \int \frac{d^2k_\perp}{ (2\pi)^3}
	\frac{dk^+}{ 2k^+}\sum_\sigma \Big[ b_\sigma(k) u({k,\sigma})
	\nonumber \\
	&  \times e^{-i(k^+\xi^--\vec{k}_\perp\cdot\vec{\xi}_\perp)}
	+ d_\sigma^\dagger(k) v({k,\sigma})e^{i(k^+\xi^--\vec{k}_\perp\cdot\vec{\xi}_\perp)}
	\Big],
\end{align}
其中 $b^\dagger(k)$ 与 $d^\dagger(k)$（$b(k)$ 与 $d(k)$）分别是夸克与反夸克的产生（湮灭）算符。{$\sigma$ 是夸克的光锥螺旋度，可取 $+1/2$ 或 $-1/2$。}粒子态及产生湮灭算符采用协变归一化，即
\begin{align}
	& \big\{b_\sigma(k),b_{\sigma'}^\dagger(k')\big\}=
	\big\{d_\sigma(k),d_{\sigma'}^\dagger(k')\big\} \nonumber \\
	& =(2\pi)^3\delta_{\sigma\sigma'}2k^+\delta(k^+-k'^{+})
	\delta^{(2)}(\vec{k}_\perp-\vec{k}_\perp^\prime)\,.
\end{align}
将上述展开代入 \eq{LCqPDF}，可得夸克分布
\begin{align}
	q(x) = \frac{1}{2x} \sum_\sigma \int \frac{d^2{\vec k}_\perp}{(2\pi)^3}
	\langle P|b_\sigma^\dagger(x,\vec{k}_\perp)b_\sigma(x,\vec{k}_\perp)|P\rangle/\langle P|P\rangle
	\label{eq:partonfock}
\end{align}
（$x>0$）；$x<0$ 时类似地得到反夸克分布。因子 $1/x$ 来自产生湮灭算符的归一化。上面的矩阵元应理解为波包态中的矩阵元在确定动量态极限下的结果~\cite{Collins:2011zzd}。这样，我们就在 LF 关联（算符）形式中恢复了 PDF 的物理含义。

\subsection{部分子结构的其他途径}
\label{sec:intro_other}

从 QCD 出发计算强子的部分子结构一直是强子物理的重要目标。除各种唯象学与模型外，主要有两条途径：光前量子化与格点 QCD。这里作者对本综述主题之外的 LFQ 与格点途径作一非常简短的回顾。

虽然 LFQ 显式使用部分子自由度，但在实际计算中并不很成功。首先，LF 微扰论如同标准的哈密顿微扰论一样明显破坏 Lorentz 对称性，需要精细的重正化方案来恢复它。一个可行的重正化方案必须处理 $\xi^-$ 方向的长程关联，这要求重正化抵消项具有泛函依赖~\cite{Wilson:1994fk}。因此 LF 微扰论从未被用于一圈以上的计算，唯一例外是 QED 中两圈反常磁矩~\cite{Langnau:1992zj}。事实上，对横向积分用维数正规化（DR）、对纵向积分用截断正规化的通行做法，虽已被 Mueller 成功用于从夸偶素波函数导出 BFKL 演化~\cite{Mueller:1993rr}，但至今未证明对多圈计算有用。

在 QCD 中使用 LFQ 的热情不在于微扰论，而在于求解强子态。文献~\cite{Pauli:1985ps}提出离散化的 LFQ 以对束缚态问题进行实际计算。这一非微扰方法对 1+1 维模型相当成功，例如 Schwinger 模型~\cite{McCartor:1994im,Harada:1995au}、1+1 维 QCD~\cite{Burkardt:1989wy,Srivastava:2000cf}、1+1 维 $\phi^4$ 理论~\cite{Harindranath:1987db}与 sine-Gordon 模型~\cite{Burkardt:1992sz}。对 3+1 维理论，曾考虑过简单近似，如 Tamm-Dancoff 近似~\cite{Perry:1990mz}。而对 QCD 本身，仍不得不对 Fock 态数目作严厉截断。此类较近的工作包括~\cite{Vary:2009gt,Lan:2019vui,Jia:2019iyk}。然而，导出完全重正化的哈密顿量十分困难，况且迄今尚无从证明 Fock 空间截断确实收敛~\cite{Wilson:1994fk}。因此，LFQ 中 QCD 束缚态的系统近似方法仍有待发现。

鉴于格点 QCD 的快速发展，用它计算部分子物理是很自然的想法。然而直接模拟实时演化在数值上是困难的，会遇到所谓的符号问题或更一般的 NP-hard 问题。此前已有若干间接计算 PDF 的方法被提出，包括已被充分研究的矩方法、强子张量与 Compton 振幅方法、坐标空间因子化等。这些途径计算的格点可观测量可通过 OPE 或色散关系与 PDF/结构函数联系起来，因而可用于探查强子部分子结构的某些信息。但其目标主要是获取 PDF 的低阶矩和/或某些坐标关联的片段，而非直接以部分子自由度表述。

格点上最常采用的途径是把 PDF 各阶矩作为局域算符矩阵元来计算~\cite{Kronfeld:1984zv,Martinelli:1987si}。在矩方法中，出发点是所谓扭度二（twist-two）算符~\cite{Christ:1972ms}
\begin{align}
	O^{\mu_1...\mu_n}= {\overline \psi}\gamma^{(\mu_1}iD^{\mu_2} ... iD^{\mu_n)} \psi
	- {\rm trace}
	\label{eq:t2opq}
\end{align}
（夸克情形），其中 $(\mu_1...\mu_2)$ 表示所有指标对称化，迹项指含至少一个度规张量因子 $g^{\mu_i \mu_j}$ 且乘以带 $n-2$ 个 Lorentz 指标的 $(n+2)$ 维算符之类的项。它们在质子态中的矩阵元为
\begin{align}
	\langle P|O^{\mu_1...\mu_n}(\mu)|P\rangle =  2a_n(\mu) (P^{\mu_1}\cdots P^{\mu_n} \!-\! {\rm trace}) \,,
	\label{eq:t2mom}
\end{align}
而 PDF 通过
\begin{align}
	{a_n(\mu)} &= { \int^1_{-1} dx x^{n-1} q(x, \mu^2)}\nn\\
	&= \int^1_0 x^{n-1} \big[q(x, \mu^2)+(-1)^n \bar q(x, \mu^2)\big]
	\label{eq:t2sum}
\end{align}
与局域矩阵元相联系，$n=1,2,...$。把 \eq{t2mom} 中所有分量都取为 $+$ 分量，
\begin{align}
	\langle P|O^{+...+}(\mu)|P\rangle =2 a_n(\mu) P^+\cdots P^+ \,,
\end{align}
并把各阶矩打包成一个分布，即可恢复 Eq. (\ref{eq:LCqPDF}) 中 PDF 的时间依赖关联。类似地，对胶子 PDF，其矩同样与局域算符的矩阵元相关：
\begin{align}
	O_g^{\mu_1...\mu_n} = -F^{(\mu_1\alpha}iD^{\mu_2}\cdots .iD^{\mu_{n-1}}F^{\mu_n)}_{\ \alpha}\ ,
	\label{eq:t2opg}
\end{align}
$n=2,4,6,...$。

迄今为止已有大量 PDF 矩的格点 QCD 计算，系统误差的控制程度不一~\cite{Lin:2017snn}，包括离散化误差、物理 $\pi$ 介子质量、有限体积效应、激发态污染以及恰当的重正化。大多数格点计算集中于非极化分布的一、二阶矩 $\langle x\rangle$~\cite{Green:2012ud,Alexandrou:2017oeh,Bali:2014gha} 与 $\langle x^2\rangle$\cite{Dolgov:2002zm,Deka:2008xr}，以及极化分布的零阶与一阶矩 $\langle 1\rangle$~\cite{Gong:2015iir,Alexandrou:2017oeh,Chang:2018uxx,Alexandrou:2019brg} 与 $\langle x\rangle$~\cite{Aoki:2010xg,Abdel-Rehim:2015owa}。但由于幂次发散与信号快速衰减，更高阶矩难以计算。尽管如此，矩计算可为任何综合性的格点 PDF 方案提供有用的校准。

为获得 PDF 更多的信息，有人提出在欧氏空间计算 DIS 的强子张量，再把结果解析延拓到 Minkowski 空间~\cite{Liu:1993cv,Liu:1998um,
	Liu:1999ak,Liu:2016djw,Liu:2017lpe,Liu:2020okp}。由于解析延拓的数值方法难以精度控制（类似前面提到的 NP-hard 或符号问题），该途径主要用于核子低激发态。以此方式获取部分子物理极具挑战性。

文献~\cite{Aglietti:1998ur,Martinelli:1998hz}提出了一个类似方法，称为``无 OPE 的算符乘积展开（OPE）''，另见~\cite{Dawson:1997ic,Capitani:1998fe,Capitani:1999fm}。要点在于非色散区的 Compton 振幅可以在欧氏空间中计算~\cite{Ji:2001wha}。通过色散关系以及在 $\nu = P\cdot q=0$ 处的 Taylor 展开，可以从格点 Compton 振幅提取结构函数的高阶矩。这一途径的最新工作与文献见~\cite{Chambers:2017dov,Hannaford-Gunn:2020pvu,Horsley:2020ltc}。类似的方法也被用于含轻重流的 Compton 振幅~\cite{Detmold:2005gg}。该途径还被用来计算 $\pi$ 介子分布振幅的二阶矩~\cite{Detmold:2018kwu,Detmold:2020lev}。

流—流关联函数也可以通过坐标空间中的 OPE 来研究，无需向流中插入动量~\cite{Braun:2007wv}。小距离处的空间关联可用于计算介子分布振幅的高阶矩。相关工作见~\cite{Braun:2015axa,Bali:2018spj,Bali:2019dqc}。更晚近，Qiu 等人~\cite{Ma:2017pxb}对部分子分布提出了类似策略，并已用于格点模拟~\cite{Sufian:2019bol,Sufian:2020vzb}。pseudo-PDF 基于 LaMET 中使用的等时关联（或 Fourier 空间的准 PDF）~\cite{Radyushkin:2017cyf,Radyushkin:2019owq}，并如~\cite{Braun:2007wv}那样使用小距离的坐标空间因子化或 OPE。鉴于其与准 PDF 的密切联系，我们将在 \sec{renorm-coord} 讨论 pseudo-PDF 数据分析方法与准 PDF 方法的比较。

关于格点上``准''夸克 TMDPDF 各阶矩也已有开创性研究~\cite{Hagler:2009mb,Musch:2010ka,Musch:2011er,Engelhardt:2015xja,Yoon:2017qzo}。钉书形（staple-shaped）规范链算符被用来连接沿空间方向分开的夸克场，以模拟 TMDPDF 的矩。这些矩的比值被假定与未知的软函数无关，可与实验数据比较。然而，在 LaMET 之前，这些构造与物理 TMDPDF 各矩之间的严格关系尚未被考察，特别是大动量极限与快度截断之间的关系——后者是 TMD 物理的关键要素。该途径与 LaMET 的比较将在 \sec{quasitmd} 给出。
%------------------------------------------------------------------------------
